\subsection{模的交换子与双交换子}

设 E 是一个环，B 是 E 的一个子集。B 在 E 中的交换子（或中心化子）是 E 中由与 B 的每个元素交换的元素组成的子环 \(B'\)。把 \(B''\)（即 \(B'\) 的交换子）称为 B 的双交换子（或双中心化子）。我们有 \(B \subset B''\)，且 \(B'\) 与其双交换子重合（III, §1, No.~2, 第 429 页）。E 的子环 B 的中心是 \(B \cap B'\)；\(B'\) 与 \(B''\) 的公共中心是 \(B' \cap B''\)。若 B 是 E 的交换子环，则我们有 \(B \subset B'\)，且 \(B''\) 是 \(B'\) 的中心（III, §1, No.~2, 第 430 页）。

设 A 是一个环，M 是左（相应地，右）A-模。我们把这些定义应用于下述情形：E 是 M 的加群的自同态环，B 是 M 的位似环 \(A_{M}\)。M 的交换子是指 \(A_{M}^{\prime}\)，即 \(A_{M}\) 在 E 中的交换子；它是 A-模 M 的自同态环。M 的双交换子是指 \(A_{M}^{\prime\prime}\)，即 \(A_{M}\) 在 E 中的双交换子；它是 M 的反模的自同态环（VIII, 第 8 页，定义 3）。我们有 \(A_{M} \subset A_{M}^{\prime\prime}\)，环 \(A_{M} \cap A_{M}^{\prime}\) 是 \(A_{M}\) 的中心，而 \(A_{M}^{\prime} \cap A_{M}^{\prime\prime}\) 是 \(A_{M}^{\prime}\) 与 \(A_{M}^{\prime\prime}\) 的公共中心。

定义 1. —— 若我们有 \(\mathrm{A}_{\mathrm{M}} = \mathrm{A}_{\mathrm{M}}^{\prime \prime}\)，则称 A-模 M 是平衡的。

若 A-模 M 是平衡的，则环 \(A_{M}\) 与 \(A_{M}^{\prime}\) 有相同的中心 \(A_{M} \cap A_{M}^{\prime}\)。模 M 是平衡 A-模当且仅当它是平衡 \(A_{M}\)-模。A-模 M 的反模是忠实的且平衡的。

当环 A 交换时，M 的双交换子 \(A_{M}^{\prime\prime}\) 是 \(\mathrm{A}_{\mathrm{M}}^{\prime}=\mathrm{End}_{\mathrm{A}}(\mathrm{M})\) 的中心；M 是平衡的意思是 \(\mathrm{End}_{\mathrm{A}}(\mathrm{M})\) 的中心化为位似。

对 A 的每个元素 \(a\)，用 \(\delta_{a}\) 表示从 A 到 A 的右位似 \(x\mapsto xa\)，用 \(\gamma_{a}\) 表示左位似 \(x\mapsto ax\)（I, §8, No.~1, 第 97 页）。映射 \(a\mapsto \delta_{a}\) 是从 \(\mathrm{A}^{\circ}\) 到左 A-模 \(\mathrm{A}_s\) 的交换子的环同构（参见 II, §10, No.~7, 第 349 页应用于 \(\mathrm{I} = \{1\}\)）。映射 \(a\mapsto \gamma_{a}\) 是从 A 到右 A-模 \(\mathrm{A}_d\) 的交换子的环同构（同前所引）。若经此映射把 A 与 \(\mathrm{A}_d\) 的交换子等同起来，则 \(\mathrm{A}_d\) 的反模与 \(\mathrm{A}_s\) 等同；因此，A-模 \(\mathrm{A}_d\) 是平衡的。同样，A-模 \(\mathrm{A}_s\) 也是平衡的。

设 n 是整数 \(\geqslant 1\)。我们把 \(A^{n}\) 视为左 \(\mathbf{M}_{n}(\mathrm{A})\)-模（同前所引）。把 \(m \in \mathbf{M}_{n}(\mathrm{A})\) 映到 \(x \mapsto mx\)（A-模 \(A_{d}^{n}\) 的自同态）的映射是从 \(\mathbf{M}_{n}(\mathrm{A})\) 到右 A-模 \(A_{d}^{n}\) 的交换子的环同构（同前所引）。

命题 1. —— 设 \((\mathrm{A}_{i})_{i\in\mathrm{I}}\) 是一族环，且对每个 \(i\in I\)，设 \(M_{i}\) 是一个 \(A_{i}\)-模。令 \(A=\prod A_{i}, M=\prod M_{i}\)，\(N=\bigoplus M_{i}\)。给 M 赋以作用律为 \(((a_{i}),(x_{i}))\mapsto(a_{i}x_{i})\) 的 A-模结构。集合 N 是 M 的 A-子模。

\begin{enumerate}
\def\labelenumi{\alph{enumi})}
\item
  映射 \((u_i) \mapsto \prod u_i\)（从 \(\prod \mathrm{End}_{\mathbf{Z}}(\mathrm{M}_i)\) 到 \(\mathrm{End}_{\mathbf{Z}}(\mathrm{M})\)，II, §1, No.~5, 第 200 页）分别限制为从 \(\prod (\mathrm{A}_i)_{\mathrm{M}_i}\)、\(\prod (\mathrm{A}_i)'_{\mathrm{M}_i}\) 与 \(\prod (\mathrm{A}_i)'''_{\mathrm{M}_i}\) 到 \(\mathrm{A}_{\mathrm{M}}\)、\(\mathrm{A}_{\mathrm{M}}'\) 与 \(\mathrm{A}_{\mathrm{M}}''\) 的环同构。
\item
  映射 \((u_i) \mapsto \bigoplus u_i\)（从 \(\prod \operatorname{End}_{\mathbf{Z}}(\mathrm{M}_i)\) 到 \(\operatorname{End}_{\mathbf{Z}}(\mathrm{N})\)，II, §1, No.~6, 第 203 页）分别限制为从 \(\prod (\mathrm{A}_i)_{\mathrm{M}_i}\)、\(\prod (\mathrm{A}_i)'_{\mathrm{M}_i}\) 与 \(\prod (\mathrm{A}_i)'''_{\mathrm{M}_i}\) 到 \(\mathrm{A_N}\)、\(\mathrm{A_N}'\) 与 \(\mathrm{A_N}''\) 的环同构。
\end{enumerate}

命题 1 中定义的环同构称为典范同构。经这些同构，乘积 \(\prod (\mathrm{A}_i)_{\mathrm{M}_i}\) 常与 \(\mathrm{A_M}\) 等同，\(\prod (\mathrm{A}_i)'_{\mathrm{M}_i}\) 与 \(\mathrm{A_M}'\) 等同，等等。

映射 \(\varphi : (u_i) \mapsto \prod u_i\)（从 \(\prod \operatorname{End}_{\mathbf{Z}}(\mathrm{M}_i)\) 到 \(\operatorname{End}_{\mathbf{Z}}(\mathrm{M})\)）是单的环同态。由 M 上 A-模结构的定义，我们有 \(\varphi(\prod (\mathrm{A}_i)_{\mathrm{M}_i}) = \mathrm{A}_{\mathrm{M}}\)。设 \(u \in \mathrm{A}_{\mathrm{M}}'\)。对每个 \(i \in \mathrm{I}\)，用 \(h_i\) 表示 A 中除指标 \(i\) 的分量为 0 外其余分量均为 1 的元素。若 \(x\) 是 M 的一个其指标 \(i\) 的分量为 0 的元素，则我们有 \(x = h_i x\)，从而 \(\operatorname{pr}_i(u(x)) = \operatorname{pr}_i(u(h_i x)) = \operatorname{pr}_i(h_i u(x)) = 0\)。因此，存在唯一的群同态 \(u_i: \mathrm{M}_i \mapsto \mathrm{M}_i\) 使 \(\operatorname{pr}_i(u(y)) = u_i(\operatorname{pr}_i(y))\) 对每个 \(y \in \mathrm{M}\) 成立。我们有 \(u = \prod u_i\)。由于映射 \(u\) 是 A-线性的，故映射 \(u_i\) 是 \(\mathrm{A}_i\)-线性的（对每个 \(i \in \mathrm{I}\)）。这就证明 \(\prod (\mathrm{A}_i)'_{\mathrm{M}_i}\) 在 \(\varphi\) 下的像包含 \(\mathrm{A}_{\mathrm{M}}'\)；反向包含是明显的。把这一结果应用于 M 的反模，便推出 \(\varphi\) 诱导出从 \(A_{M}^{\prime \prime}\) 到 \(\prod_i(A_i)_{M_i}^{\prime \prime}\) 的同构。这就证明了断言 a)。

b) 的证明与 a) 的证明经必要修改后相同。

命题 2. —— 设 A 是一个环，M 是一个 A-模。设 I 是一个集合。则 A-模 \(\mathrm{M}^{(I)}\) 的双交换子与 \(A_{M}^{\prime\prime}\)-模 \(\mathrm{M}^{(I)}\) 的位似环重合。

对 \(i \in I\)，把从 \(M^{(I)}\) 到 M 的投影同态记作 \(\pi_{i} : (x_{j})_{j \in I} \mapsto x_{i}\)，把相应的典范单射（I, §4, No.~9, 第 47 页）记作 \(\iota_{i} : M \to M^{(I)}\)。

设 \(u\) 是 \(\mathrm{End}_{\mathrm{A}}(\mathrm{M}^{(\mathrm{I})})\) 的一个元素。对所有 \(i,j \in \mathrm{I}\)，复合 \(u_{i,j} = \pi_j \circ u \circ \iota_i\) 属于交换子 \(\mathrm{A}_{\mathrm{M}}'\)（\(\mathrm{M}\) 的）。对每个元素 \(b\)（属于 \(\mathrm{A}_{\mathrm{M}}''\)）及每个 \((x_i) \in \mathrm{M}^{(\mathrm{I})}\)，我们有关系式

\[
b u \left(\left(x _ {i}\right) _ {i \in \mathrm{I}}\right) = b \left(\sum_ {i \in \mathrm{I}} u _ {i, j} \left(x _ {i}\right)\right) _ {j \in \mathrm{I}} = \left(\sum_ {i \in \mathrm{I}} u _ {i, j} \left(b x _ {i}\right)\right) _ {j \in \mathrm{I}} = u \left(b \left(x _ {i}\right) _ {i \in \mathrm{I}}\right).
\]

因此位似 \(b_{\mathrm{M}^{(\mathrm{I})}}\) 属于 A-模 \(\mathrm{M}^{(\mathrm{I})}\) 的双交换子。

反之，设 b 是 \(\mathrm{M}^{(\mathrm{I})}\) 的双交换子的一个元素。对 \(i,j\in I\)，令 \(b_{i,j}=\pi_{j}\circ b\circ\iota_{i}\)。取 \(i,j\in I\) 且 \(i\neq j\)。由于 \(\iota_{j}\circ\pi_{j}\) 属于 A-模 \(\mathrm{M}^{(\mathrm{I})}\) 的交换子，我们有

\[
b _ {i, j} = \pi_ {j} \circ \iota_ {j} \circ \pi_ {j} \circ b \circ \iota_ {i} = \pi_ {j} \circ b \circ \iota_ {j} \circ \pi_ {j} \circ \iota_ {i} = 0.
\]

同样，我们有

\[
b _ {j, j} = \pi_ {j} \circ b \circ \iota_ {j} = \pi_ {i} \circ \iota_ {i} \circ \pi_ {j} \circ b \circ \iota_ {j} = \pi_ {i} \circ b \circ \iota_ {i} \circ \pi_ {j} \circ \iota_ {j} = b _ {i, i}.
\]

此外，\(b_{i,i}\) 属于 \(A_{M}^{\prime\prime}\)。由此推出，b 与 \(A_{M}^{\prime\prime}\)-模 \(M^{(I)}\) 的一个位似重合。

